\section{教学目标与核心素养}

基于建构主义教学理念与认知负荷理论\citep{李家清2015}，本案例设定三维教学目标：

\textbf{知识目标}：掌握GIS空间分析基本原理（缓冲区分析、叠加分析、网络分析、适宜性评价）；理解生成式AI在地理数据处理与空间建模中的应用逻辑；掌握多因子加权评价模型的构建方法。

\textbf{能力目标}：能运用AI工具辅助完成城市选址类项目的完整GIS分析流程；具备利用AI生成空间分析代码、解决实际问题的能力；能对AI生成结果进行批判性评估与修正优化；能够撰写规范的GIS分析报告。

\textbf{素养目标}：培养"人机协同"思维，提升批判性使用AI工具的意识；强化地理实践力与综合思维；形成严谨的科研态度与空间思维能力；为未来从事地理信息相关科研工作奠定基础。

核心素养对接方面，本案例紧扣地理科学专业培养要求，注重培养学生的人地协调观、区域认知能力和地理实践力\citep{教育部2022}。